\section{Operators on Inner Product Spaces}

\subsection{Self-Adjoint and Normal Operators}

\begin{exercise}{7a1}
    Suppose $n$ is a positive integer.
    Define $T \in \lin{\fbb^n}$ by
    \[
        T(z_1, \ldots, z_n) = (0, z_1, \ldots, z_{n-1}).
    \]
    Find a formula for $T^\ast(z_1, \ldots, z_n)$.
\end{exercise}

\begin{sol}
    If $n = 1$, then $T(z_1) = 0$, so $T^\ast(z_1) = 0$.
    Let $n > 1$.
    Then for all $(z_1, \ldots, z_n), (w_1, \ldots, w_n) \in \fbb^n$, we have
    \[
        \gen{T(z_1, \ldots, z_n), (w_1, \ldots, w_n)} = \gen{(0, z_1, \ldots, z_{n-1}), (w_1, \ldots, w_n)} = z_1\bar{w_2} + \cdots + z_{n-1}\bar{w_n}.
    \]
    We need
    \[
        \gen{(z_1, \ldots, z_n), T^\ast(w_1, \ldots, w_n)} = z_1\bar{w_2} + \cdots + z_{n-1}\bar{w_n}.
    \]
    By comparison, we see that $T^\ast(w_1, \ldots, w_n) = (w_2, \ldots, w_n, 0)$, so $T^\ast(z_1, \ldots, z_n) = (z_2, \ldots, z_n, 0)$ for all $(z_1, \ldots, z_n) \in \fbb^n$.
\end{sol}

\begin{exercise}{7a2}
    Suppose $T \in \lin{V, W}$.
    Prove that
    \[
        T = 0 \iff T^\ast = 0 \iff T^\ast T = 0 \iff TT^\ast = 0.
    \]
\end{exercise}

\begin{sol}
    Since $(T^\ast)^\ast = T$, then $T = 0 \iff T^\ast = 0$.
    Moreover, if $T^\ast = 0$, then $T^\ast T = 0$.
    Conversely, if $T^\ast T = 0$, then for all $v \in V$, we have
    \[
        \norm{Tv}^2 = \gen{Tv, Tv} = \gen{v, T^\ast Tv} = \gen{v, 0} = 0.
    \]
    Hence, $Tv = 0$ for all $v \in V$, so $T = 0$.
    Therefore,
    \[
        T = 0 \iff T^\ast = 0 \iff T^\ast T = 0.
    \]
    Replacing $T$ with $T^\ast$ in the equivalence $T = 0 \iff T^\ast T = 0$ gives $T^\ast = 0 \iff TT^\ast = 0$.
    We then have
    \[
        T = 0 \iff T^\ast = 0 \iff T^\ast T = 0 \iff TT^\ast = 0. \qedhere
    \]
\end{sol}

\begin{exercise}{7a3}
    Suppose $T \in \lin{V}$ and $\lambda \in \fbb$.
    Prove that
    \[
        \lambda \text{ is an eigenvalue of } T \iff \bar{\lambda} \text{ is an eigenvalue of } T^\ast.
    \]
\end{exercise}

\begin{sol}
    Since $V$ is finite-dimensional,
    \[
        \dim \nul(T - \lambda I) = \dim V - \dim \range(T - \lambda I) = \dim(\range(T - \lambda I)^\perp) = \dim \nul((T - \lambda I)^\ast).
    \]
    Hence,
    \[
        \nul(T - \lambda I) \neq \set 0 \iff \nul((T - \lambda I)^\ast) \neq \set 0.
    \]
    Since $(T - \lambda I)^\ast = T^\ast - \bar{\lambda} I$, we have
    \[
        \lambda \text{ is an eigenvalue of } T \iff \bar{\lambda} \text{ is an eigenvalue of } T^\ast. \qedhere
    \]
\end{sol}

\begin{exercise}{7a4}
    Suppose $T \in \lin{V}$ and $U$ is a subspace of $V$.
    Prove that
    \[
        U \text{ is invariant under } T \iff U^\perp \text{ is invariant under } T^\ast.
    \]
\end{exercise}

\begin{sol}
    \begin{align*}
        U^\perp \text{ is invariant under } T^\ast &\iff T^\ast v \in U^\perp \text{ for all } v \in U^\perp \\
        &\iff \gen{u, T^\ast v} = 0 \text{ for all } u \in U, v \in U^\perp \\
        &\iff \gen{Tu, v} = 0 \text{ for all } u \in U, v \in U^\perp \\
        &\iff Tu \in (U^\perp)^\perp \text{ for all } u \in U \\
        &\iff Tu \in U \text{ for all } u \in U \\
        &\iff U \text{ is invariant under } T
    \end{align*}
\end{sol}

\begin{exercise}{7a5}
    Suppose $T \in \lin{V, W}$.
    Suppose $e_1, \ldots, e_n$ is an orthonormal basis of $V$ and $f_1, \ldots, f_m$ is an orthonormal basis of $W$.
    Prove that
    \[
        \norm{Te_1}^2 + \cdots + \norm{Te_n}^2 = \norm{T^\ast f_1}^2 + \cdots + \norm{T^\ast f_m}^2.
    \]
    \note{The numbers $\norm{Te_1}^2, \ldots, \norm{Te_n}^2$ in the equation above depend on the orthonormal basis $e_1, \ldots, e_n$, but the right side of the equation does not depend on $e_1, \ldots, e_n$.
    Thus, the equation above shows that the sum on the left side does not depend on which orthonormal basis $e_1, \ldots, e_n$ is used.}

    \remark{For students familiar with this topic already, note that the left-hand side is the squared Frobenius norm of the matrix representation of $T$ with respect to the orthonormal bases $e_1, \ldots, e_n$ and $f_1, \ldots, f_m$.
    This exercise shows that $\sum_i \norm{Te_i}^2 = \sum_j \norm{T^\ast f_j}^2$, which is independent of the choice of orthonormal basis $e_1, \ldots, e_n$.
    Moreover, noting that $\tr(T^\ast T) = \tr(TT^\ast)$ by the cyclic property of the trace, we see that the sum of the squared norms of the columns of the matrix representation of $T$ equals the sum of the squared norms of the rows, which aligns with the result of this exercise.}
\end{exercise}

\begin{sol}
    Let
    \[
        Te_j = \sum_{k=1}^m \gen{Te_j, f_k} f_k \implies \norm{Te_j}^2 = \sum_{k=1}^m \abs{\gen{Te_j, f_k}}^2.
    \]
    Similarly, we may write
    \[
        T^\ast f_i = \sum_{\ell=1}^n \gen{T^\ast f_i, e_\ell} e_\ell \implies \norm{T^\ast f_i}^2 = \sum_{\ell=1}^n \abs{\gen{T^\ast f_i, e_\ell}}^2.
    \]
    Note that
    \[
        \gen{T^\ast f_i, e_j} = \gen{f_i, Te_j} = \bar{\gen{Te_j, f_i}} \implies \abs{\gen{T^\ast f_i, e_j}}^2 = \abs{\gen{Te_j, f_i}}^2.
    \]
    Therefore, we have
    \[
        \sum_{k=1}^m \norm{T^\ast f_k}^2 = \sum_{k=1}^m \sum_{\ell=1}^n \abs{\gen{Te_\ell, f_k}}^2 = \sum_{\ell=1}^n \sum_{k=1}^m \abs{\gen{Te_\ell, f_k}}^2 = \sum_{\ell=1}^n \norm{Te_\ell}^2. \qedhere
    \]
\end{sol}

\begin{exercise}{7a6}
    Suppose $T \in \lin{V, W}$.
    Prove that
    \begin{subexercises}
        \item $T$ is injective $\iff$ $T^\ast$ is surjective;
        \item $T$ is surjective $\iff$ $T^\ast$ is injective.
    \end{subexercises}
\end{exercise}

\begin{solenum}
    \item $T$ injective $\iff \nul T = \set 0 \iff (\range T^\ast)^\perp = \set 0 \iff \range T^\ast = V \iff T^\ast$ surjective.
    \item Using part (a) on $T^\ast \in \lin{W, V}$ and the fact that $(T^\ast)^\ast = T$, we have
    \[
        T^\ast \text{ is injective } \iff (T^\ast)^\ast \text{ is surjective } \iff T \text{ is surjective}.
    \]
\end{solenum}

\begin{exercise}{7a7}
    Prove that if $T \in \lin{V, W}$, then
    \begin{subexercises}
        \item $\dim \nul T^\ast = \dim \nul T + \dim W - \dim V$;
        \item $\dim \range T^\ast = \dim \range T$.
    \end{subexercises}
\end{exercise}

\begin{solenum}
    \item By Theorem 7.6(a), we have $\nul T^\ast = (\range T)^\perp$, so $\dim \nul T^\ast = \dim W - \dim \range T$.
    Using the fundamental theorem of linear maps, we have $\dim \range T = \dim V - \dim \nul T$, so
    \[
        \dim \nul T^\ast = \dim W - \dim V + \dim \nul T.
    \]
    \item We have $\range T^\ast = (\nul T)^\perp$, so $\dim \range T^\ast = \dim V - \dim \nul T = \dim \range T$.
\end{solenum}

\begin{exercise}{7a8}
    Suppose $A$ is an $m$-by-$n$ matrix with entries in $\fbb$.
    Use (b) in Exercise 7 to prove that the row rank of $A$ equals the column rank of $A$.

    \note{This exercise asks for yet another alternative proof of a result that was previously proved in 3.57 and 3.133.}
\end{exercise}

\begin{sol}
    Let $T \in \lin{\fbb^n, \fbb^m}$ have the matrix $A = \mcal(T)$ with respect to the standard bases of $\fbb^n$ and $\fbb^m$.
    Since the standard bases of $\fbb^n$ and $\fbb^m$ are orthonormal, $A^\ast = \mcal(T^\ast)$ are the conjugate-transposed rows of $A$.
    Since complex conjugation preserves linear independence, the column rank of $A^\ast$ equals the row rank of $A$.
    By \exref{7a7}(b), $\dim \range T^\ast = \dim \range T$, so the column rank of $A^\ast$ equals the column rank of $A$, and the result follows.
\end{sol}

\begin{exercise}{7a9}
    Prove that the product of two self-adjoint operators on $V$ is self-adjoint if and only if the two operators commute.
\end{exercise}

\begin{sol}
    Let $S, T \in \lin{V}$ be self-adjoint, so $S^\ast = S$ and $T^\ast = T$.
    Note that $(ST)^\ast = T^\ast S^\ast = TS$.
    Hence,
    \[
        ST \text{ is self-adjoint } \iff (ST)^\ast = ST \iff TS = ST. \qedhere
    \]
\end{sol}

\begin{exercise}{7a10}
    Suppose $\fbb = \cbb$ and $T \in \lin{V}$.
    Prove that $T$ is self-adjoint if and only if
    \[
        \gen{Tv, v} = \gen{T^\ast v, v}
    \]
    for all $v \in V$.
\end{exercise}

\begin{soliff}
    \leftimp If $T$ is self-adjoint, then $T = T^\ast$, so for all $v \in V$,
    \[
        \gen{Tv, v} = \gen{T^\ast v, v}.
    \]
    \rightimp Suppose $\gen{Tv, v} = \gen{T^\ast v, v}$ for all $v \in V$.
    Then for all $v \in V$,
    \[
        \gen{(T - T^\ast)v, v} = 0.
    \]
    Since $\fbb = \cbb$, we may use 7.13 with $T - T^\ast$ in place of $T$ to conclude that $T - T^\ast = 0$, so $T = T^\ast$, i.e., $T$ is self-adjoint.
\end{soliff}

\begin{exercise}{7a11}
    Define an operator $S : \fbb^2 \to \fbb^2$ by $S(w, z) = (-z, w)$.
    \begin{subexercises}
        \item Find a formula for $S^\ast$.
        \item Show that $S$ is normal but not self-adjoint.
        \item Find all eigenvalues of $S$.
    \end{subexercises}
    \note{If $\fbb = \rbb$, then $S$ is the operator on $\rbb^2$ of counterclockwise rotation by $90^\circ$.}
\end{exercise}

\begin{solenum}
    \item For any $(w, z), (x, y) \in \fbb^2$,
    \[
        \gen{S(w, z), (x, y)} = \gen{(-z, w), (x, y)} = -z\bar x + w\bar y.
    \]
    Suppose $S^\ast(x, y) = (u, v)$. Then
    \[
        \gen{S(w, z), (x, y)} = \gen{(w, z), S^\ast(x, y)} = \gen{(w, z), (u, v)} = w\bar u + z\bar v.
    \]
    Comparing with $-z\bar x + w\bar y$, we get $w\bar u + z\bar v = -z\bar x + w\bar y$ for all $w, z \in \fbb$.
    Hence, $u = y$ and $v = -x$, so $S^\ast(x, y) = (y, -x)$.
    \item Observe that
    \[
        SS^\ast(w, z) = S(S^\ast(w, z)) = S(z, -w) = (w, z),
    \]
    and
    \[
        S^\ast S(w, z) = S^\ast(S(w, z)) = S^\ast(-z, w) = (w, z).
    \]
    Hence, $SS^\ast = S^\ast S$, so $S$ is normal.
    However,
    \[
        S(1, 0) = (0, 1) \neq (0, -1) = S^\ast(1, 0).
    \]
    Hence, $S$ is not self-adjoint.
    \item If $S(w, z) = \lambda (w, z)$ for some $\lambda \in \fbb$ and $(w, z) \neq (0, 0)$, then $-z = \lambda w$ and $w = \lambda z$.
    From this, we get $w = -\lambda^2 w$, so $\lambda^2 = -1$.
    If $\fbb = \rbb$, then $S$ has no real eigenvalues.
    If $\fbb = \cbb$, then the eigenvalues of $S$ are $\lambda = i$ and $\lambda = -i$.
\end{solenum}

\begin{exercise}{7a12}
    An operator $B \in \lin{V}$ is called \term{skew} if
    \[
        B^\ast = -B.
    \]
    Suppose that $T \in \lin{V}$.
    Prove that $T$ is normal if and only if there exist commuting operators $A$ and $B$ such that $A$ is self-adjoint, $B$ is a skew operator, and $T = A + B$.
\end{exercise}

\begin{soliff}
    \leftimp If $T$ is normal, then $TT^\ast = T^\ast T$.
    Consider the operators
    \[
        A = \frac{T + T^\ast}{2}, \quad B = \frac{T - T^\ast}{2}.
    \]
    Then $A$ is self-adjoint, $B$ is skew, and $T = A + B$.
    Moreover,
    \begin{align*}
        AB &= \frac{T + T^\ast}{2} \cdot \frac{T - T^\ast}{2} \\
        &= \frac{1}{4}(TT - TT^\ast + T^\ast T - T^\ast T^\ast) \\
        &= \frac{1}{4}(TT + TT^\ast - T^\ast T - T^\ast T^\ast) \\
        &= \frac{T - T^\ast}{2} \cdot \frac{T + T^\ast}{2} = BA,
    \end{align*}
    hence $A$ and $B$ commute.

    \rightimp Suppose there exist commuting operators $A$ and $B$ such that $A$ is self-adjoint, $B$ is skew, and $T = A + B$.
    Then
    \[
        T^\ast = (A + B)^\ast = A^\ast + B^\ast = A - B.
    \]
    Hence,
    \[
        TT^\ast = (A + B)(A - B) = A^2 - AB + BA - B^2 = A^2 - B^2,
    \]
    and
    \[
        T^\ast T = (A - B)(A + B) = A^2 + AB - BA - B^2 = A^2 - B^2.
    \]
    Therefore, $TT^\ast = T^\ast T$, so $T$ is normal.
\end{soliff}

\begin{exercise}{7a13}
    Suppose $\fbb = \rbb$.
    Define $\acal \in \lin{\lin{V}}$ by $\acal T = T^\ast$ for all $T \in \lin{V}$.
    \begin{subexercises}
        \item Find all eigenvalues of $\acal$.
        \item Find the minimal polynomial of $\acal$.
    \end{subexercises}
\end{exercise}

\begin{solenum}
    \item If $V = \set 0$, then $\lin{V} = \set 0$ and $\acal$ has no eigenvalues.
    Otherwise, assume $\dim V \geq 1$.
    If $\acal T = \lambda T$ for some $T \in \lin{V}$ and $\lambda \in \rbb$, then $T^\ast = \lambda T$.
    Taking the adjoint of both sides gives
    \[
        T = \lambda T^\ast = \lambda (\lambda T) = \lambda^2 T.
    \]
    Hence, $\lambda^2 = 1$, so $\lambda = 1$ or $\lambda = -1$.
    If $\dim V = 1$, then $T = \lambda I$ for some $\lambda \in \rbb$.
    In this case, $T^\ast = T$, so $\acal T = T = 1 \cdot T$, and the only eigenvalue of $\acal$ is $1$.
    If $\dim V \geq 2$, let $e_1, \ldots, e_n$ be an orthonormal basis of $V$ with $n = \dim V$.
    Define $T \in \lin{V}$ by $T e_1 = e_2$, $T e_2 = -e_1$, and $T e_j = 0$ for $j \geq 3$.
    Then $T^\ast e_1 = -e_2$, $T^\ast e_2 = e_1$, and $T^\ast e_j = 0$ for $j \geq 3$, so $T^\ast = -T$.
    Hence, $-1$ is also an eigenvalue of $\acal$.
    Therefore, the eigenvalues of $\acal$ are $1$ and $-1$.

    \item Observe that $\acal(\acal T) = \acal T^\ast = (T^\ast)^\ast = T$ for all $T \in \lin{V}$.
    Hence, $\acal^2 - I = 0$, so the minimal polynomial of $\acal$ must divide $z^2 - 1$.
    If $\dim V = 1$, then $\acal$ has only one eigenvalue, $1$, so the minimal polynomial of $\acal$ is $z - 1$.
    Otherwise, both eigenvalues $1$ and $-1$ occur when $\dim V \geq 2$, so the minimal polynomial of $\acal$ is $z^2 - 1$.
\end{solenum}

\begin{exercise}{7a14}
    Define an inner product on $\pcal_2(\rbb)$ by $\gen{p, q} = \int_0^1 pq$.
    Define an operator $T \in \lin{\pcal_2(\rbb)}$ by
    \[
        T(ax^2 + bx + c) = bx.
    \]
    \begin{subexercises}
        \item Show that with this inner product, the operator $T$ is not self-adjoint.
        \item The matrix of $T$ with respect to the basis $1, x, x^2$ is
        \[
            \begin{pmatrix}
                0 & 0 & 0 \\
                0 & 1 & 0 \\
                0 & 0 & 0
            \end{pmatrix}.
        \]
        This matrix equals its conjugate transpose, even though $T$ is not self-adjoint.
        Explain why this is not a contradiction.
    \end{subexercises}
\end{exercise}

\begin{solenum}
    \item Consider $p = 1$ and $q = x$. Then
    \[
        \gen{Tp, q} = \gen{0, x} = 0,
    \]
    but
    \[
        \gen{p, Tq} = \gen{1, x} = \int_0^1 x \dd x = \frac{1}{2} \neq 0.
    \]
    Hence, $T$ is not self-adjoint.
    \item Observe that
    \[
        \norm{x}^2 = \gen{x, x} = \int_0^1 x^2 \dd x = \frac{1}{3}.
    \]
    By 7.9, given an orthonormal basis, $T$ is self-adjoint if and only if the matrix of $T$ with respect to that basis equals its conjugate transpose.
    However, the basis $1, x, x^2$ is not orthonormal with respect to the given inner product, so the fact that the matrix of $T$ with respect to this basis equals its conjugate transpose does not imply that $T$ is self-adjoint.
\end{solenum}

\begin{exercise}{7a15}
    Suppose $T \in \lin{V}$ is invertible.
    Prove that
    \begin{subexercises}
        \item $T$ is self-adjoint $\iff$ $T\inv$ is self-adjoint;
        \item $T$ is normal $\iff$ $T\inv$ is normal.
    \end{subexercises}
\end{exercise}

\begin{sol}
    We first prove that $(T\inv)^\ast = (T^\ast)\inv$.
    For any $u, v \in V$, then
    \[
        \gen{u, v} = \gen{T\inv Tu, v} = \gen{Tu, (T\inv)^\ast v},
    \]
    and
    \[
        \gen{u, v} = \gen{u, T^\ast (T^\ast)\inv v} = \gen{Tu, (T^\ast)\inv v}.
    \]
    Since $\range T = V$, equality in the second components implies that $(T\inv)^\ast = (T^\ast)\inv$.
    \begin{subexercises}
        \item
        \begin{align*}
            T\inv \text{ is self-adjoint} &\iff (T\inv)^\ast = T\inv \\
            &\iff (T^\ast)\inv = T\inv \\
            &\iff T^\ast = T \\
            &\iff T \text{ is self-adjoint}.
        \end{align*}
        \item
        \begin{align*}
            T\inv \text{ is normal} &\iff (T\inv)^\ast T\inv = T\inv (T\inv)^\ast \\
            &\iff (T^\ast)\inv T\inv = T\inv (T^\ast)\inv \\
            &\iff (T^\ast T)\inv = (TT^\ast)\inv \\
            &\iff T^\ast T = TT^\ast \\
            &\iff T \text{ is normal}. \qedhere
        \end{align*}
    \end{subexercises}
\end{sol}

\begin{exercise}{7a16}
    Suppose $\fbb = \rbb$.
    \begin{subexercises}
        \item Show that the set of self-adjoint operators on $V$ is a subspace of $\lin{V}$.
        \item What is the dimension of the subspace of $\lin{V}$ in (a) [in terms of $\dim V$]?
    \end{subexercises}
\end{exercise}

\begin{solenum}
    \item Let $\scal$ be the set of self-adjoint operators on $V$.
    Since $0^\ast = 0$, we have $0 \in \scal$.
    For $S, T \in \scal$ and $\lambda \in \rbb$, we have
    \[
        (S + T)^\ast = S^\ast + T^\ast = S + T \in \scal,
    \]
    and
    \[
        (\lambda T)^\ast = \bar{\lambda} T^\ast = \lambda T \in \scal.
    \]
    Therefore, $\scal$ is closed under addition and scalar multiplication, and contains the zero operator, so it is a subspace of $\lin{V}$.
    \item Let $n = \dim V$, and let $e_1, \dots, e_n$ be an orthonormal basis of $V$.
    By 7.9, it follows that $T$ is self-adjoint if and only if $\mcal(T)$ is a symmetric matrix, i.e., $\mcal(T) = \mcal(T)^t$.
    Since a symmetric $n \times n$ matrix is determined by the entries on and above the main diagonal, the dimension of the subspace of symmetric matrices is
    \[
        \frac{n(n+1)}{2}. \qedhere
    \]
\end{solenum}

\begin{exercise}{7a17}
    Suppose $\fbb = \cbb$.
    Show that the set of self-adjoint operators on $V$ is not a subspace of $\lin{V}$.
\end{exercise}

\begin{sol}
    Note the result is not true when $V = \set 0$, since the only operator on $V$ is the zero operator, which is self-adjoint.
    Let $\scal$ be the set of self-adjoint operators on $V$.
    If $\scal$ were a subspace of $\lin{V}$, consider $iI$, where $I \in \scal$ is the identity operator on $V$.
    Then
    \[
        (iI)^\ast = \bar\imath I = -iI \neq iI.
    \]
    Therefore, $iI$ is not self-adjoint, which shows that the set of self-adjoint operators on $V$ is not a subspace of $\lin{V}$ when $V \neq \set 0$.
\end{sol}

\begin{exercise}{7a18}
    Suppose $\dim V \geq 2$.
    Show that the set of normal operators on $V$ is not a subspace of $\lin{V}$.
\end{exercise}

\begin{sol}
    Fix an orthonormal pair $e_1, e_2 \in V$.
    Extend this to an orthonormal basis $e_1, \ldots, e_n$ of $V$.
    Consider the operators $S$ and $T$ defined by
    \[
        Se_1 = e_2, Se_2 = e_1, Se_i = 0 \text{ for } i > 2, \quad Te_1 = -e_2, Te_2 = e_1, Te_i = 0 \text{ for } i > 2.
    \]
    Then
    \[
        \mcal(S) = 
        \begin{pmatrix}
            0 & 1 & \cdots & 0 \\
            1 & 0 & \cdots & 0 \\
            \vdots & \vdots & \ddots & \vdots \\
            0 & 0 & \cdots & 0
        \end{pmatrix}.
    \]
    From this, we see that $\mcal(S) = \mcal(S)^t$, so $S$ is normal.
    Moreover,
    \[
        \mcal(T) = 
        \begin{pmatrix}
            0 & 1 & \cdots & 0 \\
            -1 & 0 & \cdots & 0 \\
            \vdots & \vdots & \ddots & \vdots \\
            0 & 0 & \cdots & 0
        \end{pmatrix}.
    \]
    Direct computation yields $\mcal(T)\mcal(T)^\ast = \mcal(T)^\ast\mcal(T)$, so $T$ is normal.
    However, $S + T$ is not normal, as can be seen by computing
    \[
        \mcal(S + T) = 
        \begin{pmatrix}
            0 & 2 & \cdots & 0 \\
            0 & 0 & \cdots & 0 \\
            \vdots & \vdots & \ddots & \vdots \\
            0 & 0 & \cdots & 0
        \end{pmatrix}, \quad
        \mcal(S + T)^\ast = 
        \begin{pmatrix}
            0 & 0 & \cdots & 0 \\
            2 & 0 & \cdots & 0 \\
            \vdots & \vdots & \ddots & \vdots \\
            0 & 0 & \cdots & 0
        \end{pmatrix}.
    \]
    Direct computation shows that $\mcal(S + T)\mcal(S + T)^\ast \neq \mcal(S + T)^\ast\mcal(S + T)$, so $S + T$ is not normal.
\end{sol}

\begin{exercise}{7a19}
    Suppose $T \in \lin{V}$ and $\norm{T^\ast v} \leq \norm{Tv}$ for every $v \in V$.
    Prove that $T$ is normal.

    \note{This exercise fails on infinite-dimensional inner product spaces, leading to what are called \term{hyponormal operators}, which have a well-developed theory.}
\end{exercise}

\begin{sol}
    If $\dim V = 0$, then $T$ is trivially normal. So assume $\dim V > 0$.
    Note that for all $v \in V$, we have
    \[
        \norm{T^\ast v}^2 = \gen{T^\ast v, T^\ast v} = \gen{TT^\ast v, v}, \quad \norm{Tv}^2 = \gen{Tv, Tv} = \gen{T^\ast Tv, v}.
    \]
    From the hypothesis, we have $\norm{T^\ast v}^2 \leq \norm{Tv}^2$ for all $v \in V$, which gives
    \[
        \gen{(T^\ast T - TT^\ast)v, v} \leq 0 \quad \text{for all } v \in V.
    \]
    Let $S = TT^\ast - T^\ast T$.
    Then $S^\ast = (TT^\ast - T^\ast T)^\ast = TT^\ast - T^\ast T = S$, so $S$ is self-adjoint.
    Moreover, $\gen{Sv, v} \geq 0$ for all $v \in V$.
    Let $e_1, \ldots, e_n$ be an orthonormal basis of $V$.
    From \exref{7a5}, then
    \[
        \sum_{i=1}^n \norm{Te_i}^2 = \sum_{i=1}^n \gen{T^\ast Te_i, e_i} = \sum_{i=1}^n \gen{TT^\ast e_i, e_i} = \sum_{i=1}^n \norm{T^\ast e_i}^2.
    \]
    Then
    \[
        \sum_{i=1}^n \gen{Se_i, e_i} = \sum_{i=1}^n \gen{(TT^\ast - T^\ast T)e_i, e_i} = \sum_{i=1}^n \gen{TT^\ast e_i, e_i} - \sum_{i=1}^n \gen{T^\ast Te_i, e_i} = 0.
    \]
    Since each $\gen{Se_i, e_i} \geq 0$ and their sum is $0$, it follows that $\gen{Se_i, e_i} = 0$ for all $i$.
    For arbitrary $u \in V$ with $\norm u = 1$, we may extend it to an orthonormal basis $u, f_2, \ldots, f_n$ of $V$.
    Then
    \[
        \gen{Su, u} + \sum_{i=2}^n \gen{Sf_i, f_i} = \sum_{i=1}^n \gen{Se_i, e_i} = 0.
    \]
    Since each term is nonnegative, it follows that $\gen{Su, u} = 0$ for all unit vectors $u \in V$.
    By homogeneity, $\gen{Sv, v} = 0$ for all $v \in V$.
    From 7.16, if $S$ is self-adjoint and $\gen{Sv, v} = 0$ for all $v \in V$, then $S = 0$. 
    Hence, $TT^\ast - T^\ast T = 0$, so $TT^\ast = T^\ast T$, and therefore $T$ is normal.
\end{sol}

\begin{exercise}{7a20}
    Suppose $P \in \lin{V}$ is such that $P^2 = P$.
    Prove that the following are equivalent.
    \begin{subexercises}
        \item $P$ is self-adjoint.
        \item $P$ is normal.
        \item There is a subspace $U$ of $V$ such that $P = P_U$.
    \end{subexercises}
\end{exercise}

\begin{sol}
    \leavevmode

    \noindent (a) $\implies$ (b):
    If $P$ is self-adjoint, then $P = P^\ast$.
    Hence,
    \[
        PP^\ast = PP = P^2 = P = P^\ast P = P^\ast P,
    \]
    so $P$ is normal.

    \noindent (b) $\implies$ (c):
    Suppose $P$ is normal.
    By \exref{5a8}, $P$ is diagonalizable with eigenvalues $0$ and $1$.
    Let $U = E(1, P)$, the eigenspace corresponding to the eigenvalue $1$.
    To see that $U = \range P$, note that if $u \in U$, then $Pu = u$, so $u \in \range P$.
    Conversely, if $v \in \range P$, then $v = Pw$ for some $w \in V$.
    Then $Pv = P(Pw) = P^2 w = Pw = v$, so $v \in U$.
    Hence, $U = \range P$.
    By \exref{6c9}, $\nul P$ and $U$ are orthogonal.
    Then $\nul P \subseteq U^\perp$, and by definition of orthogonal complement, $\dim U^\perp = \dim V - \dim U$.
    Since $\dim \nul P = \dim V - \dim \range P = \dim V - \dim U = \dim U^\perp$, it follows that $\nul P = U^\perp$.
    Then for $v \in V$, we may write $v = u + w$ with $u \in U$ and $w \in U^\perp$.
    Hence,
    \[
        Pv = P(u + w) = Pu + Pw = u + 0 = u \in U.
    \]
    Therefore, $P = P_U$.

    \noindent (c) $\implies$ (a):
    Write $V = U \oplus U^\perp$.
    Then for $v = u_1 + w_1$ and $w = u_2 + w_2$ with $u_1, u_2 \in U$ and $w_1, w_2 \in U^\perp$, we have
    \[
        \gen{P_U v, w} = \gen{u_1, u_2 + w_2} = \gen{u_1, u_2} = \gen{u_1 + w_1, u_2} = \gen{v, P_U w}.
    \]
    Hence, $P_U$ is self-adjoint.
\end{sol}

\begin{exercise}{7a21}
    Suppose $D : \pcal_8(\rbb) \to \pcal_8(\rbb)$ is the differentiation operator defined by $Dp = p'$.
    Prove that there does not exist an inner product on $\pcal_8(\rbb)$ that makes $D$ a normal operator.
\end{exercise}

\begin{sol}
    If there did exist an inner product on $\pcal_8(\rbb)$ that makes $D$ a normal operator, then $\norm{Dp} = \norm{D^\ast p}$ for all $p \in \pcal_8(\rbb)$, by 7.20.
    Consider $p = x$ so that $Dp = 1 \neq 0$, but $D^2 p = 0$.
    Then applying 7.20 to $Dp$, we get
    \[
        \norm{D^\ast (Dp)} = \norm{D(Dp)} = \norm{D^2 p} = \norm{0} = 0
    \]
    so that $D^\ast (Dp) = 0$.
    Then
    \[
        \gen{D^\ast (Dp), p} = \gen{Dp, Dp} = \norm{Dp}^2.
    \]
    Since $D^\ast (Dp) = 0$, the left-hand side of the above equation is $0$, while the right-hand side is $\norm{Dp}^2 \neq 0$.
    This is a contradiction, so no such inner product can exist.
\end{sol}

\begin{exercise}{7a22}
    Give an example of an operator $T \in \lin{\rbb^3}$ such that $T$ is normal but not self-adjoint.
\end{exercise}

\begin{sol}
    Let $T$ have the matrix
    \[
        \begin{pmatrix}
            0 & -1 & 0 \\
            1 & 0 & 0 \\
            0 & 0 & 1
        \end{pmatrix}
    \]
    with respect to the standard basis of $\rbb^3$, which is orthonormal.
    Then $T$ is normal because
    \[
        T^\ast = 
        \begin{pmatrix}
            0 & 1 & 0 \\
            -1 & 0 & 0 \\
            0 & 0 & 1
        \end{pmatrix}
        = \mcal(T)^\ast.
    \]
    By direct computation, we have
    \[
        TT^\ast = I, \quad T^\ast T = I.
    \]
    Hence, $T$ is normal but not self-adjoint, since $T \neq T^\ast$.
\end{sol}

\begin{exercise}{7a23}
    Suppose $T$ is a normal operator on $V$.
    Suppose also that $v, w \in V$ satisfy the equations
    \[
        \norm{v} = \norm{w} = 2, \quad Tv = 3v, \quad Tw = 4w.
    \]
    Show that $\norm{T(v + w)} = 10$.
\end{exercise}

\begin{sol}
    Since $v, w$ are eigenvectors of the normal operator $T$ corresponding to distinct eigenvalues $3$ and $4$, they must be orthogonal.
    Therefore, $\gen{v, w} = 0$.
    Then
    \begin{align*}
        \norm{T(v + w)}^2 &= \gen{Tv + Tw, Tv + Tw} \\
        &= \gen{3v + 4w, 3v + 4w} \\
        &= \gen{3v, 3v} + \gen{3v, 4w} + \gen{4w, 3v} + \gen{4w, 4w} \\
        &= 9\gen{v, v} + 12\gen{v, w} + 12\gen{w, v} + 16\gen{w, w} \\
        &= 9\norm{v}^2 + 16\norm{w}^2 \\
        &= 9 \cdot 4 + 16 \cdot 4 \\
        &= 100.
    \end{align*}
    Therefore, $\norm{T(v + w)} = 10$.
\end{sol}

\begin{exercise}{7a24}
    Suppose $T \in \lin{V}$ and
    \[
        a_0 + a_1 z + a_2 z^2 + \cdots + a_{m-1} z^{m-1} + z^m
    \]
    is the minimal polynomial of $T$.
    Prove that the minimal polynomial of $T^\ast$ is
    \[
        \bar{a_0} + \bar{a_1} z + \bar{a_2} z^2 + \cdots + \bar{a_{m-1}} z^{m-1} + z^m.
    \]
    \note{This exercise shows that the minimal polynomial of $T^\ast$ equals the minimal polynomial of $T$ if $\fbb = \rbb$.}
\end{exercise}

\begin{sol}
    Define $p(z)$ to be the minimal polynomial of $T$.
    Let $q(z) = \sum_{i=0}^m \bar{a_i} z^i$.
    Then
    \[
        q(T^\ast) = \sum_{i=0}^m \bar{a_i} (T^\ast)^i = \sum_{i=0}^m \bar{a_i} (T^i)^\ast = \left( \sum_{i=0}^m a_i T^i \right)^\ast = p(T)^\ast = 0.
    \]
    Hence, $q(z)$ is a polynomial that annihilates $T^\ast$.
    Using a similar argument with $r(z)$ being the minimal polynomial of $T^\ast$, we can show that $\bar r(T) = 0$.
    Since $q(z)$ annihilates $T^\ast$ and $r(z)$ is the minimal polynomial of $T^\ast$, it follows that $r(z)$ divides $q(z)$.
    Similarly, since $\bar r(T) = 0$, $p(z)$ divides $\bar r(z)$.
    Then $\deg r \leq \deg q$ from the fact that $r(z)$ divides $q(z)$, and $\deg q \leq \deg r$ from the fact that $p(z)$ divides $\bar r(z)$ and $\deg p = \deg q$.
    Hence, $\deg r = \deg q$, and since $r(z)$ is monic and divides $q(z)$, we must have $r(z) = q(z)$.
    Therefore, the minimal polynomial of $T^\ast$ is $q(z)$.
\end{sol}

\begin{exercise}{7a25}
    Suppose $T \in \lin{V}$.
    Prove that $T$ is diagonalizable if and only if $T^\ast$ is diagonalizable.
\end{exercise}

\begin{soliff}
    \rightimp Suppose $T$ is diagonalizable.
    Then the minimal polynomial of $T$ splits into distinct linear factors by 5.62.
    By \exref{7a24}, the minimal polynomial of $T^\ast$ also splits into distinct linear factors whose roots are the complex conjugates of the roots of the minimal polynomial of $T$.
    Hence, $T^\ast$ is diagonalizable.

    \leftimp Conversely, suppose $T^\ast$ is diagonalizable.
    Apply the same argument as above, using the fact that $(T^\ast)^\ast = T$.
    Therefore, $T$ is diagonalizable.
\end{soliff}

\begin{exercise}{7a26}
    Fix $u, x \in V$.
    Define $T \in \lin{V}$ by $Tv = \gen{v, u}x$ for every $v \in V$.
    \begin{subexercises}
        \item Prove that if $V$ is a real vector space, then $T$ is self-adjoint if and only if the list $u, x$ is linearly dependent.
        \item Prove that $T$ is normal if and only if the list $u, x$ is linearly dependent.
    \end{subexercises}
\end{exercise}

\begin{sol}
    Note that
    \[
        \gen{Tv, w} = \gen{\gen{v, u}x, w} = \gen{v, u}\gen{x, w} = \gen{v, \bar{\gen{x, w}}u} = \gen{v, T^\ast w}.
    \]
    Since $\gen{v, u_1} = \gen{v, u_2}$ for all $v \in V$ implies $u_1 = u_2$, we conclude that $T^\ast w = \bar{\gen{x, w}}u = \gen{w, x}u$.
    \begin{subexercises}
        \item
        \rightimp If $T$ is self-adjoint, then $T = T^\ast$, which means $\gen{v, u}x = \gen{v, x}u$ for all $v \in V$.
        If $u = 0$ or $x = 0$, then the list $u, x$ is linearly dependent.
        If both are nonzero, then letting $v = u$ gives $\norm{u}^2 x = \gen{u, x}u$, hence
        \[
            x = \frac{\gen{u, x}}{\norm{u}^2} u.
        \]
        Hence, $x$ and $u$ are linearly dependent.

        \leftimp Conversely, suppose one of $u$ or $x$ is zero. Then $T = 0 = T^\ast$, so $T$ is self-adjoint.
        If both are nonzero and linearly dependent, let $x = \lambda u$ for some $\lambda \in \rbb$. Then for all $v \in V$,
        \[
            Tv = \gen{v, u}x = \gen{v, u} \lambda u = \lambda \gen{v, u} u = T^\ast v,
        \]
        hence $T$ is self-adjoint.
        \item
        \rightimp Suppose $T$ is normal. Then for any $y \in V$,
        \[
            TT^\ast y = T(\gen{y, x}u) = \gen{y, x}Tu = \gen{y, x}\gen{u, u}x = \norm{u}^2 \gen{y, x}x,
        \]
        and
        \[
            T^\ast T y = T^\ast(\gen{y, u}x) = \gen{y, u}T^\ast x = \gen{y, u}\gen{x, x}u = \norm{x}^2 \gen{y, u}u.
        \]
        If $u = 0$ or $x = 0$, then the list $u, x$ is linearly dependent.
        Otherwise, for nonzero $u$ and $x$, we must have $\norm{u}^2 \gen{y, x}x = \norm{x}^2 \gen{y, u}u$ for all $y \in V$.
        Let $y = u$. Then $\norm{u}^2 \gen{u, x}x = \norm{x}^2 \norm{u}^2 u$, hence $\gen{u, x}x = \norm{x}^2 u$.
        If $x \neq 0$, then
        \[
            u = \frac{\gen{u, x}}{\norm{x}^2} x,
        \]
        so $u$ and $x$ are linearly dependent.

        \leftimp If one of $u$ or $x$ is zero, then $T = 0 = T^\ast$, so $T$ is normal.
        If both are nonzero and linearly dependent, let $u = \mu x$ for some $\mu \in \fbb$.
        Then for all $y \in V$,
        \[
            TT^\ast y = \norm{u}^2 \gen{y, x}x = \norm{\mu x}^2 \gen{y, x}x = \abs{\mu}^2 \norm{x}^2 \gen{y, x}x,
        \]
        and
        \[
            T^\ast T y = \norm{x}^2 \gen{y, u}u = \norm{x}^2 \gen{y, \mu x} \mu x = \abs{\mu}^2 \norm{x}^2 \gen{y, x}x.
        \]
        Hence, $TT^\ast y = T^\ast T y$ for all $y \in V$, so $T$ is normal. \qedhere
    \end{subexercises}
\end{sol}

\begin{exercise}{7a27}
    Suppose $T \in \lin{V}$ is normal.
    Prove that
    \[
        \nul T^k = \nul T \quad \text{and} \quad \range T^k = \range T
    \]
    for every positive integer $k$.
\end{exercise}

\begin{sol}
    Let $v \in \nul T$.
    Then $T^k v = 0$, so $v \in \nul T^k$.
    Hence, $\nul T \subseteq \nul T^k$.
    Conversely, let $v \in \nul T^k$.
    Since $T$ is normal, $\norm{Tv} = \norm{T^\ast v}$ by 7.20.
    Hence, $\nul T = \nul T^\ast$.
    If $T^2 v = 0$, then $T(Tv) = 0$, so $Tv \in \nul T = \nul T^\ast$.
    Therefore, $T^\ast Tv = 0$ so that $\norm{Tv}^2 = \gen{Tv, Tv} = \gen{T^\ast Tv, v} = 0$.
    Hence, $Tv = 0$, which shows that $v \in \nul T$.
    By induction on $k$, assume that $\nul T^{k-1} = \nul T$.
    If $v \in \nul T^k$, then $T v \in \nul T^{k-1} = \nul T$, so $T(Tv) = T^2 v = 0$.
    By the base case earlier, $v \in \nul T^2 = \nul T$.
    Hence, by induction, $\nul T^k = \nul T$ for all positive integers $k$.

    For the other equality, note that $v \in \range T^k$ implies that there exists $u \in V$ such that $v = T^k u$.
    Then $T(T^{k-1} u) = T^k u = v$, so $v \in \range T$.
    Hence, $\range T^k \subseteq \range T$.
    Conversely, note that
    \[
        \dim \range T^k = \dim V - \dim \nul T^k = \dim V - \dim \nul T = \dim \range T
    \]
    by the previous equality $\nul T^k = \nul T$.
    Since $\range T^k \subseteq \range T$ and they have the same dimension, we conclude that $\range T^k = \range T$.
\end{sol}

\begin{exercise}{7a28}
    Suppose $T \in \lin{V}$ is normal.
    Prove that if $\lambda \in \fbb$, then the minimal polynomial of $T$ is not a polynomial multiple of $(x - \lambda)^2$.
\end{exercise}

\begin{sol}
    If the minimal polynomial $p(x)$ of $T$ is a polynomial multiple of $(x - \lambda)^2$, then $p(x) = (x - \lambda)^2 q(x)$ for some polynomial $q(x)$.
    Then $(T - \lambda I)^2 q(T) = 0$, but $(T - \lambda I) q(T) \neq 0$ since $p(x)$ is the minimal polynomial.
    Pick $v$ such that $(T - \lambda I) q(T) v \neq 0$.
    Then $(T - \lambda I)^2 q(T) v = 0$, so $q(T) v \in \nul (T - \lambda I)^2$.
    By normality of $T$, we know that $(T - \lambda I)^\ast = T^\ast - \bar\lambda I$ which commutes with $T - \lambda I$.
    Hence, by \exref{7a27}, $\nul (T - \lambda I)^2 = \nul (T - \lambda I)$.
    Therefore, $q(T) v \in \nul (T - \lambda I)$, which implies that $(T - \lambda I) q(T) v = 0$.
    This contradicts our choice of $v$.
    Hence, the minimal polynomial of $T$ cannot be a polynomial multiple of $(x - \lambda)^2$.
\end{sol}

\begin{exercise}{7a29}
    Prove or give a counterexample: If $T \in \lin{V}$ and there is an orthonormal basis $e_1, \ldots, e_n$ of $V$ such that $\norm{Te_k} = \norm{T^\ast e_k}$ for each $k = 1, \ldots, n$, then $T$ is normal.
\end{exercise}

\begin{sol}
    Let $V = \fbb^2$, and consider the standard orthonormal basis $e_1, e_2$.
    Define $T \in \lin{V}$ by $T(e_1) = e_1 - e_2$ and $T(e_2) = e_1$.
    Then
    \[
        \mcal(T) = 
        \begin{pmatrix}
            1 & 1 \\
            -1 & 0 \\
        \end{pmatrix}, \quad
        \mcal(T^\ast) = 
        \begin{pmatrix}
            1 & -1 \\
            1 & 0 \\
        \end{pmatrix}
    \]
    so that $T^\ast e_1 = e_1 + e_2$ and $T^\ast e_2 = -e_1$.
    Moreover,
    \[
        \norm{Te_1}^2 = \norm{e_1 - e_2}^2 = 2, \quad \norm{T^\ast e_1}^2 = \norm{e_1 + e_2}^2 = 2,
    \]
    and
    \[
        \norm{Te_2}^2 = \norm{e_1}^2 = 1, \quad \norm{T^\ast e_2}^2 = \norm{-e_1}^2 = 1.
    \]
    Hence, $\norm{Te_k} = \norm{T^\ast e_k}$ for each $k = 1, 2$.
    However, $T$ is not normal because
    \[
        \mcal(TT^\ast) =
        \begin{pmatrix}
            2 & -1 \\
            -1 & 1 \\
        \end{pmatrix}, \quad
        \mcal(T^\ast T) =
        \begin{pmatrix}
            2 & 1 \\
            1 & 1
        \end{pmatrix}.
    \]
    Therefore, $TT^\ast \neq T^\ast T$.
\end{sol}

\begin{exercise}{7a30}
    Suppose that $T \in \lin{\fbb^3}$ is normal and $T(1, 1, 1) = (2, 2, 2)$.
    Suppose $(z_1, z_2, z_3) \in \nul T$.
    Prove that $z_1 + z_2 + z_3 = 0$.
\end{exercise}

\begin{sol}
    Since $T$ is normal, we have $\nul T = \nul T^\ast$ by 7.20.
    Let $z = (z_1, z_2, z_3) \in \nul T = \nul T^\ast$.
    By $T(1, 1, 1) = 2(1, 1, 1)$ and the definition of the adjoint, we have
    \[
        2\gen{(1, 1, 1), z} = \gen{T(1, 1, 1), z} = \gen{(1, 1, 1), T^\ast z} = \gen{(1, 1, 1), 0} = 0.
    \]
    Then $\gen{(1, 1, 1), z} = 0$.
    Direct calculation shows that
    \[
        \gen{(1, 1, 1), z} = \bar{z_1} + \bar{z_2} + \bar{z_3} = \bar{z_1 + z_2 + z_3}.
    \]
    Since $\gen{(1, 1, 1), z} = 0$, we have $\bar{z_1 + z_2 + z_3} = 0$, which implies $z_1 + z_2 + z_3 = 0$.
\end{sol}

\begin{exercise}{7a31}
    Fix a positive integer $n$.
    In the inner product space of continuous real-valued functions on $[-\pi, \pi]$ with inner product $\gen{f, g} = \int_{-\pi}^{\pi} fg$, let
    \[
        V = \spn(1, \cos x, \cos 2x, \ldots, \cos nx, \sin x, \sin 2x, \ldots, \sin nx).
    \]
    \begin{subexercises}
        \item Define $D \in \lin{V}$ by $Df = f'$.
        Show that $D^\ast = -D$.
        Conclude that $D$ is normal but not self-adjoint.
        \item Define $T \in \lin{V}$ by $Tf = f''$.
        Show that $T$ is self-adjoint.
    \end{subexercises}
\end{exercise}

\begin{solenum}
    \item By a similar set of arguments as in \exref{6b4}, the list above is orthogonal with respect to the given inner product.
    Note that any $f, g \in V$ is a linear combination of $1, \cos kx, \sin kx$ for $k = 1, 2, \ldots, n$, all of which are periodic with period $2\pi$.
    Then $f(\pi) = f(-\pi)$ and $g(\pi) = g(-\pi)$.
    Integration by parts gives
    \[
        \gen{Df, g} = \int_{-\pi}^{\pi} f'g = fg\,\big|_{-\pi}^{\pi} - \int_{-\pi}^{\pi} f g' = -\int_{-\pi}^{\pi} f g' = \gen{f, -Dg}.
    \]
    Then $D^\ast = -D$, which shows that
    \[
        DD^\ast = D(-D) = -D^2 \quad\text{and}\quad D^\ast D = (-D)D = -D^2.
    \]
    Hence, $DD^\ast = D^\ast D$, which shows that $D$ is normal.
    However, $D$ is not self-adjoint, since $D\cos x = -\sin x \neq \sin x = D^\ast \cos x$.
    \item Note that $T = D^2$.
    By the previous part, $D^\ast = -D$, so $T^\ast = (D^2)^\ast = (D^\ast)^2 = (-D)^2 = D^2 = T$.
    Hence, $T$ is self-adjoint.
\end{solenum}

\begin{exercise}{7a32}
    Suppose $T : V \to W$ is a linear map.
    Show that under the standard identification of $V$ with $V'$ (see 6.58) and the corresponding identification of $W$ with $W'$, the adjoint map $T^\ast : W \to V$ corresponds to the dual map $T' : W' \to V'$.
    More precisely, show that
    \[
        T'(\phi_w) = \phi_{T^\ast w}
    \]
    for all $w \in W$, where $\phi_w$ and $\phi_{T^\ast w}$ are defined as in 6.58.
\end{exercise}

\begin{sol}
    For all $v \in V$ and $w \in W$, we have
    \[
        (T'(\phi_w))v = \phi_w(Tv) = \gen{Tv, w} = \gen{v, T^\ast w} = \phi_{T^\ast w}v.
    \]
    Hence, $T'(\phi_w) = \phi_{T^\ast w}$ for all $w \in W$.
\end{sol}